\documentclass[12pt,a4paper]{article}
\usepackage[utf8]{inputenc}
\usepackage[spanish]{babel}
\usepackage{amsmath}
\usepackage{array}
\usepackage{geometry}
\usepackage{booktabs}   % Para líneas de tablas de calidad profesional
\usepackage{float}      % Para fijar la posición de las tablas normales con [H]
\usepackage{xltabular}  % ¡NUEVO! Combina tablas de ancho fijo con multipágina

\geometry{margin=1in}

% Para centrado y alineación en columnas p{}
\newcommand{\PreserveBackslash}[1]{\let\temp=\\#1\let\\=\temp}
\newcolumntype{C}[1]{>{\PreserveBackslash\centering}p{#1}}
\newcolumntype{L}[1]{>{\PreserveBackslash\raggedright}p{#1}}
\newcolumntype{R}[1]{>{\PreserveBackslash\raggedleft}p{#1}}

% Mayor espacio vertical entre las filas de las tablas
\renewcommand{\arraystretch}{1.3}

\title{Resumen de Arquitectura MIPS}
\author{}
\date{}

\begin{document}
	
	\maketitle
	
	\section*{Formatos básicos de instrucción}
	
	\begin{table}[H] 
		\centering
		\begin{tabular}{c|c|c|c|c|c|c}
			\toprule
			\textbf{Tipo} & \multicolumn{5}{c}{\textbf{Campos}} & \\
			\midrule
			R & cod oper (31-26) & rs (25-21) & rt (20-16) & rd (15-11) & desplaz (10-6) & func (5-0) \\
			I & cod oper (31-26) & rs (25-21) & rt (20-16) & \multicolumn{3}{c}{inmediato (15-0)} \\
			J & cod oper (31-26) & \multicolumn{5}{c}{dirección (25-0)} \\
			\bottomrule
		\end{tabular}
	\end{table}
	
	\section*{Lenguaje ensamblador MIPS}
	
	% TABLA ÚNICA MULTIPÁGINA (Sin entorno \begin{table})
		\small
		\begin{xltabular}{\textwidth}{>{\raggedright\arraybackslash}p{2.2cm} >{\raggedright\arraybackslash}p{2.2cm} >{\raggedright\arraybackslash}p{3cm} >{\raggedright\arraybackslash}p{3.5cm} X}
			\caption{Instrucciones MIPS.}\\
			\toprule
			\textbf{Categoría} & \textbf{Instrucción} & \textbf{Ejemplo} & \textbf{Significado} & \textbf{Comentarios} \\
			\midrule
			\endfirsthead % Esto define la cabecera de la primera página
			
			\multicolumn{5}{c}{{\bfseries \tablename\ \thetable{} -- Continuación de la página anterior}} \\
			\toprule
			\textbf{Categoría} & \textbf{Instrucción} & \textbf{Ejemplo} & \textbf{Significado} & \textbf{Comentarios} \\
			\midrule
			\endhead % Esto repite la cabecera en las siguientes páginas
			
			\bottomrule
			\endfoot % Línea inferior en páginas intermedias
			
			\bottomrule
			\endlastfoot % Línea inferior en la última página
			
			Aritmética 
			& add & \texttt{add \$s1, \$s2, \$s3} & \texttt{\$s1 = \$s2 + \$s3} & Tres operandos; datos en registros \\
			& subtract & \texttt{sub \$s1, \$s2, \$s3} & \texttt{\$s1 = \$s2 - \$s3} & Tres operandos; datos en registros \\
			& add immed. & \texttt{addi \$s1, \$s2, 100} & \texttt{\$s1 = \$s2 + 100} & Usado para sumar constantes \\
			\midrule
			Transferencia de dato 
			& load word & \texttt{lw \$s1, 100(\$s2)} & \texttt{\$s1 = Mem[\$s2+100]} & Palabra de memoria a registro \\
			& store word & \texttt{sw \$s1, 100(\$s2)} & \texttt{Mem[\$s2+100] = \$s1} & Palabra de registro a memoria \\
			& load half & \texttt{lh \$s1, 100(\$s2)} & \texttt{\$s1 = Mem[\$s2+100]} & Media palabra de memoria a registro \\
			& store half & \texttt{sh \$s1, 100(\$s2)} & \texttt{Mem[\$s2+100] = \$s1} & Media palabra de registro a memoria \\
			& load byte & \texttt{lb \$s1, 100(\$s2)} & \texttt{\$s1 = Mem[\$s2+100]} & Byte de memoria a registro \\
			& store byte & \texttt{sb \$s1, 100(\$s2)} & \texttt{Mem[\$s2+100] = \$s1} & Byte de registro a memoria \\
			& load upper immed. & \texttt{lui \$s1, 100} & \texttt{\$s1 = 100 $\cdot 2^{16}$} & Cargar constante en los 16 bits de mayor peso \\
			\midrule
			Lógica 
			& and & \texttt{and \$s1, \$s2, \$s3} & \texttt{\$s1 = \$s2 \& \$s3} & Tres registros operandos; AND bit-a-bit \\
			& or & \texttt{or \$s1, \$s2, \$s3} & \texttt{\$s1 = \$s2 | \$s3} & Tres registros operandos; OR bit-a-bit \\
			& nor & \texttt{nor \$s1, \$s2, \$s3} & \texttt{\$s1 = $\sim$(\$s2 | \$s3)} & Tres registros operandos; NOR bit-a-bit \\
			& and immed. & \texttt{andi \$s1, \$s2, 100} & \texttt{\$s1 = \$s2 \& 100} & AND bit-a-bit registro con constante \\
			& or immed. & \texttt{ori \$s1, \$s2, 100} & \texttt{\$s1 = \$s2 | 100} & OR bit-a-bit registro con constante \\
			& shift left & \texttt{sll \$s1, \$s2, 10} & \texttt{\$s1 = \$s2 $\ll$ 10} & Desplazamiento a la izquierda por constante \\
			& shift right & \texttt{srl \$s1, \$s2, 10} & \texttt{\$s1 = \$s2 $\gg$ 10} & Desplazamiento a la derecha por constante \\
			\midrule
			Salto condicional 
			& branch on eq & \texttt{beq \$s1, \$s2, L} & if (\texttt{\$s1 == \$s2}) go L & Comprueba igualdad y bifurca relativo al PC \\
			& branch not eq & \texttt{bne \$s1, \$s2, L} & if (\texttt{\$s1 != \$s2}) go L & Comprueba si no igual y bifurca relativo al PC \\
			& set on less & \texttt{slt \$s1, \$s2, \$s3} & if (\texttt{\$s2 < \$s3}) \texttt{\$s1=1} else \texttt{\$s1=0} & Compara si es menor que; usado con beq, bne \\
			& set on less imm & \texttt{slti \$s1, \$s2, 100}& if (\texttt{\$s2 < 100}) \texttt{\$s1=1} else \texttt{\$s1=0} & Compara si es menor que una constante \\
			\midrule
			Salto incond. 
			& jump & \texttt{j 2500} & go to 10000 & Salto a la dirección destino \\
			& jump reg & \texttt{jr \$ra} & go to \texttt{\$ra} & Para retorno de procedimiento \\
			& jump and link & \texttt{jal 2500} & \texttt{\$ra = PC+4}; go 10000 & Para llamada a procedimiento \\
		\end{xltabular}
		
		\section*{Operativos MIPS}
		
		\begin{table}[H]
			\centering
			\begin{tabularx}{\textwidth}{l X}
				\toprule
				\textbf{Nombre} & \textbf{Comentarios} \\
				\midrule
				32 registros & \texttt{\$s0-\$s7}, \texttt{\$t0-\$t9}, \texttt{\$zero}, \texttt{\$a0-\$a3}, \texttt{\$v0-\$v1}, \texttt{\$gp}, \texttt{\$fp}, \texttt{\$sp}, \texttt{\$ra}, \texttt{\$at} \newline Localizaciones rápidas para los datos. En MIPS, los datos deben estar en los registros para realizar operaciones aritméticas. El registro MIPS \texttt{\$zero} es siempre igual a 0. El registro \texttt{\$at} está reservado por el ensamblador para manejar constantes grandes. \\
				\bottomrule
			\end{tabularx}
		\end{table}
		
		\section*{Nombre de registro, número, uso y convenio de llamada}
		
		\begin{table}[H]
			\centering
			\small
			\begin{tabularx}{\textwidth}{l c X c}
				\toprule
				\textbf{NOMBRE} & \textbf{NÚMERO} & \textbf{USO} & \textbf{¿SE CONSERVA?} \\
				\midrule
				\texttt{\$zero} & 0 & Valor constante 0 & N/A \\
				\texttt{\$at} & 1 & Ensamblador temporal & No \\
				\texttt{\$v0-\$v1} & 2-3 & Valores de resultado de funciones y evaluación de expresiones & No \\
				\texttt{\$a0-\$a3} & 4-7 & Argumentos & No \\
				\texttt{\$t0-\$t7} & 8-15 & Temporales & No \\
				\texttt{\$s0-\$s7} & 16-23 & Temporales almacenados & Sí \\
				\texttt{\$t8-\$t9} & 24-25 & Temporales & No \\
				\texttt{\$k0-\$k1} & 26-27 & Reservados para el núcleo del Sistema Operativo & No \\
				\texttt{\$gp} & 28 & Puntero global & Sí \\
				\texttt{\$sp} & 29 & Puntero de pila & Sí \\
				\texttt{\$fp} & 30 & Puntero de marco & Sí \\
				\texttt{\$ra} & 31 & Dirección de retorno & Sí \\
				\bottomrule
			\end{tabularx}
		\end{table}
		
	\end{document}
